% ===== PRÉAMBULE CANONIQUE — à coller VERBATIM en tête de CHAQUE .tex =====
\documentclass[11pt,a4paper]{article}
\usepackage[utf8]{inputenc}\usepackage[T1]{fontenc}\usepackage{lmodern}
\usepackage[french]{babel}
\usepackage{amsmath,amssymb,mathtools}\usepackage{icomma}
\usepackage[version=4]{mhchem}          % \ce{...} : équations & formules
\usepackage{siunitx}\sisetup{locale=FR} % \qty{}{}, \si{}, \ang{}
\usepackage{chemfig}                    % \chemfig{...} : structures développées
\usepackage{xcolor}\usepackage{colortbl}
\usepackage{tikz}\usepackage{pgfplots}\pgfplotsset{compat=1.18}
\usepackage[most]{tcolorbox}\usepackage{enumitem}\usepackage{array,booktabs}
\usepackage[a4paper,margin=2.2cm]{geometry}
\usetikzlibrary{arrows.meta,calc,angles,quotes,positioning,patterns,backgrounds,shapes.geometric}
\definecolor{primary}{RGB}{222,49,99}\definecolor{primarydark}{RGB}{150,22,55}
\definecolor{defcol}{RGB}{0,114,178}\definecolor{methcol}{RGB}{52,92,148}
\definecolor{excol}{RGB}{0,135,90}\definecolor{attncol}{RGB}{200,40,55}
\tcbset{boxbase/.style={enhanced,breakable,boxrule=0.9pt,arc=2.5pt,
  left=3mm,right=3mm,top=2.3mm,bottom=2.3mm,fonttitle=\bfseries,coltitle=white}}
% --- boîtes TUTORIEL (noms EXACTS, ne pas renommer) ---
\newtcolorbox{cle}[1][]{boxbase,colback=defcol!6,colframe=defcol,
  title={\textsc{À retenir}\ifx\relax#1\relax\relax\else\textmd{ : \itshape #1}\fi}}
\newtcolorbox{methode}[1][]{boxbase,colback=methcol!7,colframe=methcol,sharp corners,
  title={\textsc{Méthode}\ifx\relax#1\relax\relax\else\textmd{ --- \itshape #1}\fi}}
\newtcolorbox{exemplebox}[1][]{boxbase,colback=excol!5,colframe=excol,
  title={\textsc{Exemple}\ifx\relax#1\relax\relax\else\textmd{ : \itshape #1}\fi}}
\newtcolorbox{piege}[1][]{boxbase,colback=attncol!6,colframe=attncol,
  title={\textsc{Piège}\ifx\relax#1\relax\relax\else\textmd{ : \itshape #1}\fi}}
\newtcolorbox{remarque}[1][]{boxbase,colback=black!4,colframe=black!45,
  title={\textsc{Remarque}\ifx\relax#1\relax\relax\else\textmd{ : \itshape #1}\fi}}
% --- boîtes EXERCICE / CORRIGÉ (noms EXACTS, auto-comptées) ---
\newtcolorbox[auto counter]{exo}[1][]{boxbase,colback=primary!4,colframe=primary,
  title={\textsc{Exercice}~\thetcbcounter\ifx\relax#1\relax\relax\else\textmd{ --- \itshape #1}\fi}}
\newtcolorbox[auto counter]{corrige}[1][]{boxbase,colback=excol!4,colframe=excol,
  title={\textsc{Corrigé}~\thetcbcounter\ifx\relax#1\relax\relax\else\textmd{ --- \itshape #1}\fi}}
\newcommand{\R}{\mathbb{R}}\newcommand{\N}{\mathbb{N}}\newcommand{\Z}{\mathbb{Z}}\newcommand{\ds}{\displaystyle}
% =========================================================================
\begin{document}
% THEME_TITLE: Entropie et spontanéité
\begin{center}
{\Large\bfseries\color{primarydark} Thème 4 — Entropie et spontanéité}\\[2pt]
{\color{primary}\small Le désordre, le sens d'évolution, les 2\textsuperscript{e} et 3\textsuperscript{e} principes}
\end{center}

\begin{cle}[entropie et désordre]
L'\textbf{entropie} $S$ mesure le \textbf{désordre} d'un système (le nombre d'arrangements possibles). Elle s'exprime en $\si{\joule\per\mol\per\kelvin}$. Plus un système est désordonné, plus $S$ est grand :
\[ \text{solide} \;<\; \text{liquide} \;<\; \text{gaz} \;<\; \text{gaz dispersé}. \]
\textbf{2\textsuperscript{e} principe} : un processus est \textbf{spontané} si le désordre \emph{total} augmente, c'est-à-dire si $\Delta S_{\text{univers}} > 0$. \textbf{3\textsuperscript{e} principe} : à $T=\qty{0}{\kelvin}$, un cristal parfait est parfaitement ordonné, donc $S = \qty{0}{\joule\per\mol\per\kelvin}$.
\end{cle}

\begin{center}
\begin{tikzpicture}[scale=1.0,dot/.style={circle,fill=defcol,inner sep=1.5pt}]
  % solide : grille ordonnée
  \foreach \x in {0,0.35,0.7} \foreach \y in {0,0.35,0.7} \node[dot] at (\x,\y){};
  \node[font=\scriptsize] at (0.35,-0.4){solide};
  % liquide : cluster légèrement désordonné
  \begin{scope}[xshift=3.0cm]
    \foreach \p in {(0,0.05),(0.4,0.1),(0.15,0.4),(0.55,0.45),(0.3,0.7),(0.7,0.7),(0.05,0.75),(0.6,0.15),(0.35,0.25)}
      \node[dot] at \p {};
    \node[font=\scriptsize] at (0.35,-0.4){liquide};
  \end{scope}
  % gaz : dispersé
  \begin{scope}[xshift=6.0cm]
    \foreach \p in {(0.1,0.7),(0.9,0.55),(0.5,0.9),(1.2,0.2),(0.3,0.1),(1.0,0.85),(0.7,0.35)}
      \node[dot] at \p {};
    \node[font=\scriptsize] at (0.55,-0.4){gaz};
  \end{scope}
  \draw[->,line width=1pt,primary] (-0.2,-0.75)--(7.2,-0.75) node[right,font=\scriptsize,primary]{$S$ croissant};
\end{tikzpicture}
\end{center}

\begin{methode}[prévoir le signe de $\Delta S$ « à vue »]
On compare le désordre entre réactifs et produits, en regardant surtout les \textbf{gaz} :
\begin{itemize}[leftmargin=1.6em,topsep=2pt,itemsep=1.5pt]
  \item \textbf{apparition} de gaz ou \textbf{augmentation} du nombre de moles de gaz $\Rightarrow \Delta S>0$ ;
  \item \textbf{disparition} de gaz ou \textbf{diminution} du nombre de moles de gaz $\Rightarrow \Delta S<0$ ;
  \item \textbf{dissolution} d'un solide $\Rightarrow \Delta S>0$ ;
  \item \textbf{fusion / vaporisation / sublimation} $\Rightarrow \Delta S>0$ (et l'inverse $\Rightarrow \Delta S<0$).
\end{itemize}
Quand le nombre de moles de gaz ne change \emph{pas}, $\Delta S$ est \textbf{faible} (proche de $0$).
\end{methode}

\begin{exemplebox}[trois cas typiques]
\begin{itemize}[leftmargin=1.6em,topsep=2pt,itemsep=1.5pt]
  \item $\ce{N2}(g) + 3\,\ce{H2}(g) \rightarrow 2\,\ce{NH3}(g)$ : $4$ mol de gaz $\to 2$ mol $\Rightarrow \Delta S<0$.
  \item $\ce{CaCO3}(s) \rightarrow \ce{CaO}(s) + \ce{CO2}(g)$ : on crée du gaz $\Rightarrow \Delta S>0$.
  \item $\ce{CO2}(s) \rightarrow \ce{CO2}(g)$ (sublimation) : solide $\to$ gaz $\Rightarrow \Delta S>0$.
\end{itemize}
\end{exemplebox}

\begin{piege}[à ne pas confondre]
\begin{itemize}[leftmargin=1.6em,topsep=2pt,itemsep=2pt]
  \item on compte surtout les moles de \textbf{gaz} : solides et liquides « pèsent » beaucoup moins dans le désordre.
  \item \textbf{spontané $\ne$ rapide} : le 2\textsuperscript{e} principe donne le \emph{sens} d'évolution, pas la \emph{vitesse} (qui relève de la cinétique).
  \item ne pas confondre $\Delta S_{\text{système}}$ et $\Delta S_{\text{univers}}$ : c'est le \emph{signe de} $\Delta S_{\text{univers}}$ qui décide de la spontanéité.
\end{itemize}
\end{piege}

\end{document}
